\section{Variational selection and the global curvature cost}
\label{sec:variational-selection-global-cost}

The recovery metrics of Section~\ref{sec:variational-geometry} give an
upper bound for the first-order cost.  We now extract additional
information from genuine near-minimality and distinguish it from
classical curvature minimization and simultaneous monopole identities.
These results do not determine the unrestricted first-order coefficient.

Use the balanced family of Theorem~\ref{thm:vg-balanced-family}:
\begin{equation}\label{webvar:family}
 M_j=X_j\#k_j\overline{\mathbb{CP}}^{\,2}\#\ell_j(S^1\times S^3),
 \quad A_j=c_1^2(X_j),\quad Q_j=k_j+4\ell_j,\quad h_j=Q_j/A_j.
\end{equation}
Here $h_j\to0$, $Q_j\to\infty$, and $U_j=1+O(h_j)$.  Write
\begin{equation}\label{webvar:infima}
 P(g)=\int|\Ric_g^\circ|^2\,d\mu_g,\qquad
 p_j=\inf_gP(g),\qquad \alpha_j=\inf_g\E(g).
\end{equation}
Ricci norms are tensor norms; curvature and Weyl norms are
Hilbert--Schmidt norms of operators on two-forms.  The recovery and
Theorem~\ref{thm:vg-first-order-reduction} give
\begin{equation}\label{webvar:reduction}
 p_j\le8\pi^2Q_j,\qquad
 \alpha_j=\frac{4p_j}{8\pi^2A_j}+O(h_j^2),\qquad
 \int|\Rm_g|^2\,d\mu_g=8\pi^2\chi(M_j)+P(g).
\end{equation}
Thus the unresolved coefficient-four assertion is
$p_j\ge8\pi^2Q_j-o(Q_j)$.  The all-class and simultaneous-field
statements use $b_+(M_j)\ge2$, so that all ordinary sign classes
are supplied without a chamber choice.  They are not assertions
about arbitrary chamber-dependent data.

\subsection{The full ordinary monopole budget}

Fix $M=X\#k\overline{\mathbb{CP}}^{\,2}\#\ell(S^1\times S^3)$
with $b_+(M)\ge2$, and suppress $j$.  Let $\mathcal C(M)$ be the
finite set of real cohomology classes of all ordinary monopole
determinants.  Following the convex-geometric comparison in
\cite{LeBrunConvex}, define
\begin{equation}\label{webvar:all-class-definitions}
 \mathfrak b^2(M)=\max_{v\in\operatorname{conv}(\mathcal C(M))}v^2,
 \qquad
 \mathfrak a^2(M)=\inf_H\max_{a\in\mathcal C(M)}(a_H^+)^2,
\end{equation}
where $H$ ranges over maximal positive subspaces of $H^2(M;\mathbb R)$.
These symbols are separate from the variational infimum $\alpha_j$.

\begin{proposition}[Full ordinary class budget]
\label{webvar:all-class-budget}
On this surgery manifold,
\begin{equation}\label{webvar:all-class-value}
 \mathfrak b^2(M)=\mathfrak a^2(M)=A.
\end{equation}
No classification of $\mathcal C(M)$ is required.
\end{proposition}
\begin{proof}
The ordinary blow-up and handle supply contains all classes
\[
 a_\epsilon=c+\sum_{i=1}^k\epsilon_iE_i,\qquad
 c^2=A,\quad c\cdot E_i=0,\quad E_i\cdot E_m=-\delta_{im},
 \quad\epsilon_i\in\{-1,1\};
\]
see Section~\ref{sec:variational-geometry} and
\cite[Proposition~3.2]{LeBrun2001}.  Their average is $c$, whence
$\mathfrak b^2(M)\ge c^2=A$.  For any positive plane $H$ and any
convex combination $v=\sum_\nu t_\nu a_\nu$,
\[
 v^2\le |v_H^+|^2
 \le\left(\sum_\nu t_\nu|(a_\nu)_H^+|\right)^2
 \le\max_{a\in\mathcal C(M)}(a_H^+)^2.
\]
Maximizing over $v$ and then taking the infimum over $H$ gives
$\mathfrak b^2(M)\le\mathfrak a^2(M)$.

The recovery metrics of Proposition~\ref{prop:vg-recovery}, coming
from \cite[Theorem~4.3]{LeBrun2001}, satisfy
$S(g_\varepsilon)\to32\pi^2A$.  Every class in $\mathcal C(M)$
has an ordinary solution on the same metric $g_\varepsilon$.
The scalar monopole inequality thus holds simultaneously for all
these classes:
\[
 \max_{a\in\mathcal C(M)}(a_{g_\varepsilon}^+)^2
       \le S(g_\varepsilon)/(32\pi^2).
\]
The Hodge plane is admissible in \eqref{webvar:all-class-definitions},
so $\mathfrak a^2(M)\le A$.  This completes the comparison.
\end{proof}

This is an application of the classical recovery and convex-budget
method.  Additional ordinary classes cannot raise either uniform
budget above $A$, although their projected squares at a specified
metric may exceed $A$.  It does not compute the full class budget
of the different Bauer--Furuta connected sums in
Section~\ref{sec:bf-budget-comparison}; there only the supplied sign
family has been completely evaluated.

The classical energy result is
\cite[Propositions~4.1--4.2 and Theorem~4.3]{LeBrun2001}
\begin{equation}\label{webvar:classical-infima}
 \inf_gR(g)=8\pi^2(A+Q),\qquad \inf_gS(g)=32\pi^2A.
\end{equation}
One cannot subtract these infima to obtain a lower bound for
$P=R-S/4$.  Explicitly, with $W=W_++W_-$, set
\[
 \Delta_S=S-32\pi^2A,\qquad
 \Delta_W=W-\left(4\pi^2A/3+12\pi^2k\right).
\]
The reference Weyl value is its value along the same recovery.
Since $P=S/12+2W-16\pi^2\chi$ and $\chi=A/3+k-2\ell$,
\begin{equation}\label{webvar:joint-energies}
 P=8\pi^2Q+\Delta_S/12+2\Delta_W,\qquad
 R=8\pi^2(A+Q)+\Delta_S/3+2\Delta_W.
\end{equation}
The classical Ricci bound gives only $\Delta_W\ge-\Delta_S/6$.
The desired first-order bound requires
$\Delta_W\ge-\Delta_S/24-o(Q)$ on true near-minimizers.
This is a reformulation, not a proof of the missing inequality.
The distinction between total-curvature and Ricci-energy infima
is already explicit in \cite[Section~11]{LeBrun2004}.  The exact
handle-only value in Section~\ref{sec:variational-geometry} remains
a previously proved special case; it does not settle the balanced
family with $k_j/Q_j\to2/9$.

\subsection{Selection with a critical derivative budget}

For a smooth symmetric covariant two-tensor $H$, define
\begin{equation}\label{webvar:test-budget}
 \mathfrak D_g(H)=\int_M
       (|\nabla^2H|_g^2+|\nabla H|_g^4)\,d\mu_g,\qquad
 \mathcal T(g,A)=
 \{H:\|H\|_{\infty,g}\le1,\ \mathfrak D_g(H)\le A\}.
\end{equation}
Both conditions are unchanged when $g$ and $H$ are multiplied by
the same positive constant.  They allow tests on shrinking regions.

\begin{lemma}[Uniform finite-variation remainder]
\label{webvar:variation-remainder}
Suppose $g$ is smooth on a closed four-manifold and
$\int|\Rm_g|^2\,d\mu_g\le K A$.  There are a dimensional $t_0>0$
and $C_K<\infty$ depending only on $K$ and the dimension such that,
for $H\in\mathcal T(g,A)$ and $|t|\le t_0$, $g_t=g+tH$ is a metric
and
\begin{equation}\label{webvar:remainder-bound}
 \left|\frac{d^m}{dt^m}P(g_t)\right|\le C_K A,\qquad m=2,3.
\end{equation}
The constants require no injectivity-radius, diameter, Sobolev,
or pointwise derivative bound.
\end{lemma}
\begin{proof}
The eigenvalues of $g^{-1}H$ lie in $[-1,1]$, so $|t|\le1/2$
gives $g/2\le g_t\le3g/2$.  Relative to the connection of $g$,
\[
 (\nabla^{g_t}-\nabla^g)^a{}_{bc}
 =\frac t2(g_t)^{ad}
       (\nabla_bH_{cd}+\nabla_cH_{bd}-\nabla_dH_{bc}).
\]
Spatial differentiation expresses the curvature difference as fixed
contractions of $t g_t^{-1}*\nabla^2H$ and
$t^2g_t^{-1}*g_t^{-1}*\nabla H*\nabla H$.
Here $*$ denotes tensor contractions with fixed coefficients.
Changing metric contractions also introduces the original $\Rm_g$.
Differentiating in $t$ up to order three adds bounded contractions
with $g_t^{-1}$ and $H$, since
$\partial_tg_t^{-1}=-g_t^{-1}Hg_t^{-1}$.
The curvature, trace-free Ricci tensor, and required $t$ derivatives
are therefore bounded in the fixed $g$ norm by
\[
 C_m(|\Rm_g|+|\nabla^2H|+|\nabla H|^2).
\]
The volume-density ratio and its first three derivatives are bounded
as well.  Differentiating the integral defining $P$ and applying
Young's inequality to all cross terms bounds both derivatives by
\[
 C\int_M(|\Rm_g|^2+|\nabla^2H|^2+|\nabla H|^4)\,d\mu_g\le C_K A.
\]
No uniform coordinate radius was used.
\end{proof}

\begin{proposition}[Critical-budget selection of near-minimizers]
\label{webvar:selection}
For any prescribed positive sequence $\varepsilon_j$, there are
smooth metrics $g_j$ on the balanced family satisfying
\begin{equation}\label{webvar:selection-energy}
 s_{g_j}=-4,\quad 0\le P(g_j)-p_j\le\varepsilon_j,\quad
 P(g_j)-p_j\to0,\quad \E(g_j)-\alpha_j=o(h_j).
\end{equation}
They can be chosen so that, uniformly for $H\in\mathcal T(g_j,A_j)$,
\begin{equation}\label{webvar:selection-variations}
 |DP_{g_j}[H]|=o(Q_j),\qquad D^2P_{g_j}[H,H]\ge-o(Q_j).
\end{equation}
The second derivative is along $g_j+tH$; no critical metric is assumed.
\end{proposition}
\begin{proof}
Choose any $\eta_j\downarrow0$ and set
\begin{equation}\label{webvar:selection-precision}
 \delta_j=\min\{\varepsilon_j,j^{-1},\eta_j^3Q_j^3/A_j^2\}>0.
\end{equation}
Choose a smooth metric with $P-p_j\le\delta_j$ and replace it by
its scalar-$-4$ Yamabe representative.  Every conformal class is
negative Yamabe, and Proposition~\ref{prop:vg-conformal-infimum} gives
\[
 P(g)-P(\widehat g)
       =\bigl(S(g)-S(\widehat g)\bigr)/12\ge0.
\]
Scale invariance of $P$ therefore leaves
$0\le P(g_j)-p_j\le\delta_j$.  Recovery gives $P(g_j)=O(Q_j)$.
The two-sided budget \eqref{eq:vg-two-sided-R}, $R=P+4V$, and
\eqref{webvar:reduction} give
\begin{equation}\label{webvar:selected-budgets}
 R(g_j)=8\pi^2A_j+O(Q_j),\quad
 V(g_j)=2\pi^2A_j+O(Q_j),\quad
 \int|\Rm_{g_j}|^2\,d\mu_{g_j}=O(A_j).
\end{equation}
For the last assertion use $\chi(M_j)=A_j/3+O(Q_j)$.
Thus Lemma~\ref{webvar:variation-remainder} applies with fixed $K$.

For any allowed $H$, put $F(t)=P(g_j+tH)$.  Global near-minimality gives
\begin{equation}\label{webvar:finite-minimality}
 F(t)\ge p_j\ge F(0)-\delta_j\qquad(|t|\le t_0).
\end{equation}
Taylor expansion at both $t$ and $-t$, with $|F''|\le C A_j$, yields
$|F'(0)|\le\delta_j/t+C A_jt$.  Take
$t=(\delta_j/A_j)^{1/2}$, eventually below $t_0$, to get
\begin{equation}\label{webvar:first-variation-bound}
 |DP_{g_j}[H]|\le C(A_j\delta_j)^{1/2}
       \le C\eta_j^{3/2}Q_j\sqrt{h_j}=o(Q_j).
\end{equation}
Adding the second-order Taylor expansions and using $|F'''|\le C A_j$
gives $F''(0)\ge-2\delta_j/t^2-C A_jt$.  Choosing
$t=(\delta_j/A_j)^{1/3}\le\eta_jh_j$ gives
\begin{equation}\label{webvar:second-variation-bound}
 D^2P_{g_j}[H,H]\ge-C A_j^{2/3}\delta_j^{1/3}\ge-C\eta_jQ_j.
\end{equation}
Every constant is independent of $H$.  Finally, at $s=-4$ one has
$\E=4P/R$, so \eqref{webvar:reduction} and
\eqref{webvar:selected-budgets} imply
\[
 0\le\E(g_j)-\alpha_j\le C(h_j^2+\delta_j/A_j)=o(h_j).
\]
The remaining precision assertions follow from
$\delta_j\le\varepsilon_j$ and $\delta_j\le j^{-1}$.
\end{proof}

At $s=-4$, the first variation of $S$ on the closed manifold is
\[
 DS_g[H]=\int\langle-2s\Ric+(s^2/2)g,H\rangle\,d\mu_g
         =8\int\langle\Ric^\circ,H\rangle\,d\mu_g.
\]
The derivative terms vanish by integration by parts since $s$ is constant.
Using $P=4W_++S/12-8\pi^2T_+$, the selection therefore gives
\begin{equation}\label{webvar:weak-equation}
 4DW_{+,g_j}[H]+\frac23\int
       \langle\Ric_{g_j}^\circ,H\rangle\,d\mu_{g_j}=o(Q_j),
       \qquad H\in\mathcal T(g_j,A_j),
\end{equation}
uniformly.  Variational notation for $DW_+$ avoids a convention-dependent
Bach sign.  This weak identity gives no residual regularity and does
not justify $H=\Ric^\circ$ without its amplitude and derivative bounds.

\begin{corollary}[Disjoint regions]\label{webvar:disjoint-regions}
For the selected $g_j$, let finitely many $H_{ji}$ have pairwise disjoint
supports, $\|H_{ji}\|_{\infty,g_j}\le1$, and
$\sum_i\mathfrak D_{g_j}(H_{ji})\le A_j$.  Then
\begin{equation}\label{webvar:disjoint-sums}
 \sum_i|DP_{g_j}[H_{ji}]|=o(Q_j),\qquad
 \sum_i(-D^2P_{g_j}[H_{ji},H_{ji}])_+=o(Q_j).
\end{equation}
\end{corollary}
\begin{proof}
Choose signs $\epsilon_i$ so that
$\epsilon_iDP[H_{ji}]=|DP[H_{ji}]|$.  Disjointness gives
$H=\sum_i\epsilon_iH_{ji}\in\mathcal T(g_j,A_j)$, since amplitude
does not increase and both integrals in $\mathfrak D$ add.
Linearity and \eqref{webvar:selection-variations} give the first claim.
For the second, sum only tensors with negative second variation.
The density defining $P$ is local in the metric's second jet, so
mixed second variations on disjoint supports vanish.  Thus the
second variation of this admissible sum is precisely the sum of
the negative values, and \eqref{webvar:selection-variations} applies.
\end{proof}

A bounded test on a region of radius $r$ with smoothly controlled
rescaled geometry has $|\nabla H|=O(r^{-1})$,
$|\nabla^2H|=O(r^{-2})$, and support volume $O(r^4)$.
Its critical budget is $O(1)$ as $r\to0$.  This admits microscopic
and degenerating regions beyond the fixed $C^4$ compact cell
families in Section~\ref{sec:signed-selection}.  Rescaled control
verifies this example's test budget; it is not an extra assumption
of Proposition~\ref{webvar:selection}.  A fixed multiple of the
budget $A_j$ is also allowed after a fixed rescaling of the tests.
An actual descent direction with this controlled budget is still needed.

\begin{proposition}[Conditional instability criterion]
\label{webvar:instability-criterion}
Suppose that for every $\epsilon>0$ some $c_\epsilon>0$ has this
property: for all large $j$, every scalar-$-4$ metric with
$P(g)\le8\pi^2(1-\epsilon)Q_j$ admits $H\in\mathcal T(g,A_j)$ such that
\begin{equation}\label{webvar:missing-instability}
 |DP_g[H]|\ge c_\epsilon Q_j
 \quad\text{or}\quad D^2P_g[H,H]\le-c_\epsilon Q_j.
\end{equation}
Then $p_j/(8\pi^2Q_j)\to1$ and $\alpha_j/h_j\to4$.
\end{proposition}
\begin{proof}
Otherwise a subsequence has $p_j\le8\pi^2(1-2\epsilon)Q_j$ for
some $\epsilon>0$.  Proposition~\ref{webvar:selection} supplies metrics
eventually satisfying the threshold in \eqref{webvar:missing-instability},
while both alternatives contradict \eqref{webvar:selection-variations}.
The lower limit is therefore at least one; recovery and
\eqref{webvar:reduction} give both limits.
\end{proof}

Hypothesis~\eqref{webvar:missing-instability} is not established here.
Selection needs no attained minimum or compactness theorem.  It concerns
the selected sequence, not all metrics with $P=O(Q)$, and supplies
neither a strong stationary equation nor an ideal geometric limit.

\subsection{A simultaneous-field identity for the exceptional cost}

Fix $s_g=-4$ on one surgery manifold with the chamber-independent
ordinary supply.  For each $\epsilon\in\{-1,1\}^k$, choose an actual
solution on this same metric with determinant class
$a_\epsilon=c+\sum_i\epsilon_iE_i$.  In the normalization of
Section~\ref{sec:single-solution-coercivity},
\begin{equation}\label{webvar:ordinary-fields}
 D_{\mathcal A_\epsilon}\Phi_\epsilon=0,\qquad
 F_{\mathcal A_\epsilon}^+=2i\sigma(\Phi_\epsilon),\qquad
 \rho_\epsilon=|\Phi_\epsilon|^2\in[0,2].
\end{equation}
Set
\[
 f_\epsilon=-iF_{\mathcal A_\epsilon},\quad
 [f_\epsilon/(2\pi)]=a_\epsilon,\quad
 K_\epsilon=\int|\nabla_{\mathcal A_\epsilon}\Phi_\epsilon|^2,\quad
 \zeta_\epsilon=\langle\Phi_\epsilon,
                    \nabla_{\mathcal A_\epsilon}\Phi_\epsilon\rangle.
\]
All integrals and norms in this subsection use $g$.  Let
$\mathbb E_\epsilon$ be the uniform sign average, and define
\begin{equation}\label{webvar:walsh-forms}
 f_I=\mathbb E_\epsilon\left[
       \left(\prod_{i\in I}\epsilon_i\right)f_\epsilon\right],
       \qquad I\subset\{1,\ldots,k\}.
\end{equation}
These average real closed curvature forms, not connections on
different determinant bundles.

\begin{lemma}[Exact exceptional contribution]\label{webvar:walsh-identity}
\begin{equation}\label{webvar:walsh-energy}
 \mathbb E_\epsilon\|f_\epsilon^-\|_2^2
 =4\pi^2k+\|f_\varnothing^-\|_2^2
                   +\sum_{I\ne\varnothing}\|f_I^+\|_2^2.
\end{equation}
\end{lemma}
\begin{proof}
Sign orthogonality gives
$[f_\varnothing/(2\pi)]=c$, $[f_{\{i\}}/(2\pi)]=E_i$, and
$[f_I]=0$ for $|I|\ge2$ in real cohomology.  Parseval gives
$\mathbb E_\epsilon\|f_\epsilon^\pm\|_2^2=\sum_I\|f_I^\pm\|_2^2$.
Chern--Weil gives $\|f_I^-\|_2^2-\|f_I^+\|_2^2=4\pi^2$
for a singleton and zero for $|I|\ge2$.
Substitution proves the identity, with no exponential factor in $k$.
\end{proof}

For each solution put
\[
 \begin{gathered}
 p_\epsilon=(a_{\epsilon,g}^+)^2,\quad
 J_\epsilon=\|f_\epsilon^+-(f_\epsilon^+)_H\|_2^2,\quad
 \Sigma_\epsilon=16V-32\pi^2p_\epsilon,\quad
 D_\epsilon=24V+18W_+-72\pi^2p_\epsilon,\\
 \psi_\epsilon=2\sqrt2\sigma(\Phi_\epsilon),\qquad
 W_{0,\epsilon}=\frac14\left(\frac{\rho_\epsilon^2}{3}\Id
                         -\psi_\epsilon\otimes\psi_\epsilon\right).
 \end{gathered}
\]
The subscript $H$ denotes harmonic projection.  Define
\begin{equation}\label{webvar:amplitude-polynomials}
 C(t)=t+1-\frac{(t+2)^2}{12},\qquad
 b(t)=2C(t)-1=\frac{-t^2+8t+2}{6}\ge\frac13
 \quad(0\le t\le2).
\end{equation}

\begin{proposition}[Simultaneous signed cost identity]
\label{webvar:simultaneous-cost}
For every choice of the actual solutions,
\begin{equation}\label{webvar:field-cost}
 P(g)-8\pi^2Q=\mathscr A(g,\boldsymbol\Phi)
                         -4\mathbb E_\epsilon K_\epsilon,
\end{equation}
where
\begin{equation}\label{webvar:positive-field-cost}
 \begin{split}
 \mathscr A={}&2\|f_\varnothing^-\|_2^2
                 +2\sum_{I\ne\varnothing}\|f_I^+\|_2^2\\
 &+\mathbb E_\epsilon\left[
 8\|\zeta_\epsilon\|_2^2+4\|W^+-W_{0,\epsilon}\|_2^2
       +\int b(\rho_\epsilon)(\rho_\epsilon-2)^2\right]\ge0.
 \end{split}
\end{equation}
All these identities are exact for the fixed metric.
\end{proposition}
\begin{proof}
Equations~\eqref{eq:ssc-exact-decomposition} and
\eqref{eq:ssc-low-slack}, for the same solution in this normalization, are
\begin{align}
 D_\epsilon/18
 &=J_\epsilon+2\|\zeta_\epsilon\|_2^2+\|W^+-W_{0,\epsilon}\|_2^2
       +\tfrac12\int C(\rho_\epsilon)(\rho_\epsilon-2)^2,
       \label{webvar:single-decomposition}\\
 \Sigma_\epsilon
 &=16K_\epsilon+4\int(\rho_\epsilon-2)^2+8J_\epsilon.
       \label{webvar:single-slack}
\end{align}
For the second equation, \eqref{eq:ssc-spinor-rescaling} gives
$K_{\rm rescaled}=(2/\sqrt6)K_\epsilon$ and therefore
$8\sqrt6K_{\rm rescaled}=16K_\epsilon$.
Since $|f_\epsilon^+|^2=\rho_\epsilon^2/2$, harmonic projection gives
\begin{equation}\label{webvar:spinor-hodge}
 \int\rho_\epsilon^2=8\pi^2p_\epsilon+2J_\epsilon.
\end{equation}
Consequently
$4W_++4V/3-\int\rho_\epsilon^2
 =2D_\epsilon/9-\Sigma_\epsilon/4-2J_\epsilon$.
Substitute \eqref{webvar:single-decomposition} and
\eqref{webvar:single-slack}.  The $J_\epsilon$ terms cancel, leaving
\begin{equation}\label{webvar:signed-single-remainder}
 \begin{split}
 4W_++4V/3-\int\rho_\epsilon^2
  ={}&8\|\zeta_\epsilon\|_2^2+4\|W^+-W_{0,\epsilon}\|_2^2\\
    &+\int b(\rho_\epsilon)(\rho_\epsilon-2)^2-4K_\epsilon.
 \end{split}
\end{equation}
Also $a_\epsilon^2=A-k$ implies
$\|f_\epsilon^+\|_2^2-\|f_\epsilon^-\|_2^2=4\pi^2(A-k)$.
Together with $P=4W_++4V/3-8\pi^2(A-Q)$, this gives
\begin{equation}\label{webvar:chern-single-cost}
 P=32\pi^2\ell+2\|f_\epsilon^-\|_2^2
                    +4W_++4V/3-\int\rho_\epsilon^2.
\end{equation}
It also reflects $c_2(S^+_{a_\epsilon})[M]=(a_\epsilon^2-T_+)/4=\ell$.
Averaging, applying \eqref{webvar:walsh-energy}, and substituting
\eqref{webvar:signed-single-remainder} proves \eqref{webvar:field-cost},
since $32\pi^2\ell+8\pi^2k=8\pi^2Q$.
Nonnegativity follows from \eqref{webvar:amplitude-polynomials}.
\end{proof}

The determinant part also equals
\begin{equation}\label{webvar:walsh-classical-comparison}
 2\|f_\varnothing^-\|_2^2+2\sum_{I\ne\varnothing}\|f_I^+\|_2^2
 =8\pi^2(\mathbb E_\epsilon p_\epsilon-A)
                                  +2\mathbb E_\epsilon J_\epsilon.
\end{equation}
Indeed, average
$\|f_\epsilon^-\|_2^2=4\pi^2(p_\epsilon-A+k)+J_\epsilon$
and use \eqref{webvar:walsh-energy}.  Thus the simultaneous organization
does not itself improve the averaged classical projection budget.

If $P\le8\pi^2(1-\delta)Q$, with $\delta>0$, then necessarily
\begin{equation}\label{webvar:necessary-kinetic-cost}
 \mathbb E_\epsilon K_\epsilon\ge2\pi^2\delta Q+\mathscr A/4.
\end{equation}
The missing sharp estimate on a selected near-minimizing sequence is
\begin{equation}\label{webvar:missing-field-inequality}
 4\mathbb E_\epsilon K_\epsilon\le\mathscr A+o(Q).
\end{equation}
Estimates with unspecified constants and errors of order $Q$ do
not prove this signed inequality.  Solution selection alone cannot
improve it: every choice at fixed $g$ satisfies
$4\mathbb E K_\epsilon-\mathscr A=8\pi^2Q-P(g)$.
Even a smooth branch of actual fields along a metric variation
would merely recover $DP$ and $D^2P$ by differentiation, without
determining the missing constant in the value.

For $k>0$, the averaged form $f_\varnothing$ cannot supply a new
ordinary monopole equation: $c\cdot E_i=0$ but $E_i^2=-1$ is odd,
so $c$ is not characteristic and cannot be a Spin$^c$ determinant
class on this blow-up.  Curvature averaging does not remove this obstruction.

For comparison, retain the signed density of
Section~\ref{sec:signed-selection}:
\begin{equation}\label{webvar:signed-density}
 \begin{gathered}
 f_g=\frac{s_g^2/12+4|W_g^+|^2-|\Ric_g^\circ|^2}{8\pi^2},
 \qquad N(g)=\int(f_g)_-\,d\mu_g,\\
 Z(g)=\int_{\{f_g\ge0\}}|\Ric_g^\circ|^2\,d\mu_g
       +\int_{\{f_g<0\}}(s_g^2/12+4|W_g^+|^2)\,d\mu_g.
 \end{gathered}
\end{equation}
The sign partition gives $P=8\pi^2N+Z$: on $\{f_g<0\}$,
$8\pi^2(f_g)_-=|\Ric_g^\circ|^2-s_g^2/12-4|W_g^+|^2$,
while the complement contributes to $Z$.
Thus \eqref{webvar:missing-field-inequality} is equivalent on the
same sequence to the unproved compensation
\begin{equation}\label{webvar:missing-geometric-inequality}
 Z(g_j)\ge8\pi^2(Q_j-N(g_j))-o(Q_j).
\end{equation}
This retains the payment by $Z$ when $N<Q$, without assuming the
stronger selection $N\ge Q-o(Q)$ for every competitive metric.
Neither exact decomposition proves the compensation.

\subsection{Ideal minimizers and whole asymptotically Euclidean cores}

By \eqref{webvar:reduction}, $P$ and total curvature have the same
minimizing sequences on a fixed closed manifold.  Anderson's classical
ideal-minimizer theorem \cite[Theorem~4.10]{Anderson2005} gives selected
minimizing sequences with critical orbifold limits.  The compact
alternative includes an energy equality with all ALE bubbles; the
open or collapsing alternative gives an inequality.  The regular
limit minimizes against compact perturbations.  Completeness and
finitely many components in the open alternative are not supplied
by that theorem \cite[Remark~4.12]{Anderson2005}.  It does not identify
the regular bulk with $X$ or turn quotient-ended bubbles into whole
Euclidean-ended cores.  Yamabe replacement preserves the minimizing
value asymptotically, but need not preserve a given ideal decomposition.

The following relative calculation states what whole cores would
supply.  AE means that all end groups are trivial; ALE permits
nontrivial finite end groups.

\begin{lemma}[Relative curvature cost with end corrections]
\label{webvar:ale-cost}
Let $N$ be a connected smooth oriented four-manifold with finitely
many ALE ends, finite-dimensional homology, and negative-definite
intersection form.  Assume finite curvature energy and decay
sufficient for the relative Gauss--Bonnet and signature formulas.
Set $e=\#\{\text{ends}\}$, $m=b_2(N)$, $r=b_1(N)$, and
$q_\infty=\sum_{a=1}^e|\Gamma_a|^{-1}$.
Fix the signature-correction convention by
\begin{equation}\label{webvar:signature-convention}
 W_N^+-W_N^-=12\pi^2(\tau(N)+\eta_\infty).
\end{equation}
With $P_N=\int_N|\Ric^\circ|^2$ and $S_N=\int_Ns^2$,
\begin{equation}\label{webvar:ale-exact-cost}
 P_N=8\pi^2m+32\pi^2r+16\pi^2(e-2+q_\infty)
           -24\pi^2\eta_\infty+4W_N^++S_N/12.
\end{equation}
For AE ends this becomes
\begin{equation}\label{webvar:ae-cost}
 P_N=8\pi^2b_2(N)+32\pi^2(b_1(N)+e-1)+4W_N^++S_N/12.
\end{equation}
No Einstein or scalar-flat assumption is imposed.
\end{lemma}
\begin{proof}
Let $K$ be the compact core.  Its boundary components are rational
homology three-spheres.  The exact sequence of $(K,\partial K)$
and Poincar\'e--Lefschetz duality over $\mathbb R$ give
\[
 0\longrightarrow H_4(K,\partial K)\cong\mathbb R
 \longrightarrow H_3(\partial K)\cong\mathbb R^e
 \longrightarrow H_3(K)
 \longrightarrow H_3(K,\partial K)\cong H^1(K)
 \longrightarrow0.
\]
Thus $b_3(N)=e-1+r$ and
\begin{equation}\label{webvar:core-topology}
 \chi(N)=2-e-2r+m,\qquad \tau(N)=-m.
\end{equation}
The intersection form is nondegenerate because the boundary is a
union of rational homology spheres, and its signature is $-m$ by
negative definiteness.  The classical relative Gauss--Bonnet formula is
\begin{equation}\label{webvar:relative-gb}
 8\pi^2(\chi(N)-q_\infty)
       =W_N^++W_N^--P_N/2+S_N/24.
\end{equation}
Use $W_N^-=W_N^++12\pi^2(m-\eta_\infty)$ from
\eqref{webvar:signature-convention} and substitute
\eqref{webvar:core-topology}.  The result is
$P_N=4W_N^++24\pi^2(m-\eta_\infty)+S_N/12
 -16\pi^2(2-e-2r+m-q_\infty)$, which is
\eqref{webvar:ale-exact-cost}.  For AE ends,
$q_\infty=e$ and $\eta_\infty=0$, proving \eqref{webvar:ae-cost}.
\end{proof}

The classical Eguchi--Hanson metric \cite{EguchiHanson1978},
oriented with negative intersection form, has
$m=1$, $r=0$, $e=1$, $\Gamma=\mathbb Z_2$, and $P_N=S_N=0$.
The signature correction of its round $\mathbb{RP}^3$ end is zero.
Formula~\eqref{webvar:ale-exact-cost} gives
$8\pi^2+16\pi^2(-1+1/2)=0$ for the topological part, agreeing with
Ricci-flatness.  The AE bound $P_N\ge8\pi^2m$ would be false here.
Terminal quotient-ended bubbles cannot replace whole cores with
their outer marking.

\begin{proposition}[Conditional sharp cost from whole AE cores]
\label{webvar:whole-ae-sharpness}
Fix $M=X\#k\overline{\mathbb{CP}}^{\,2}\#\ell(S^1\times S^3)$.
Suppose a minimizing configuration for
$\mathcal R(g)=\int|\Rm_g|^2$ consists of a smooth closed bulk
$(B,g_B)$ diffeomorphic to $X$ and finitely many connected complete
whole AE cores $N_\nu$ satisfying Lemma~\ref{webvar:ale-cost}.
Assume that their compact cores attach to the punctured bulk along
all their Euclidean end spheres and this marked gluing gives $M$.
If the curvature accounting satisfies
\begin{equation}\label{webvar:whole-ae-accounting}
 \inf_g\mathcal R(g)\ge\mathcal R(g_B)+\sum_\nu\mathcal R(N_\nu),
\end{equation}
then
\begin{equation}\label{webvar:whole-ae-result}
 \inf_gP(g)=8\pi^2(k+4\ell).
\end{equation}
An energy equality in \eqref{webvar:whole-ae-accounting} is unnecessary.
\end{proposition}
\begin{proof}
Write $m_\nu=b_2(N_\nu)$, $r_\nu=b_1(N_\nu)$, and $e_\nu$ for
the end count.  The gluing spheres have $H_1=H_2=0$, so
Mayer--Vietoris gives a direct sum on second homology.  In first
homology it adds the cycles of the gluing graph.  The graph has
one bulk vertex, one vertex per core, and $e_\nu$ edges from
core $\nu$ to the bulk, hence first Betti number
$\sum_\nu(e_\nu-1)$.  Comparison with the prescribed connected sum gives
\begin{equation}\label{webvar:marked-counts}
 k=\sum_\nu m_\nu,\qquad
 \ell=\sum_\nu(r_\nu+e_\nu-1).
\end{equation}
The curvature decomposition and \eqref{webvar:relative-gb} give
$\mathcal R(N_\nu)=8\pi^2(\chi(N_\nu)-e_\nu)+P_{N_\nu}$.
Equations~\eqref{webvar:core-topology} and \eqref{webvar:ae-cost} yield
\[
 \mathcal R(N_\nu)
 =16\pi^2(m_\nu+r_\nu+e_\nu-1)+4W_{N_\nu}^++S_{N_\nu}/12
 \ge16\pi^2(m_\nu+r_\nu+e_\nu-1).
\]
The closed bulk satisfies
$\mathcal R(g_B)=8\pi^2\chi(X)+P(g_B)\ge8\pi^2\chi(X)$ without
an Einstein assumption.  Thus
\[
 \inf_gP(g)\ge8\pi^2\chi(X)+16\pi^2(k+\ell)-8\pi^2\chi(M)
            =8\pi^2(k+4\ell),
\]
because $\chi(M)=\chi(X)+k-2\ell$.  Recovery gives the reverse bound.
\end{proof}

For the varying balanced family the same proof gives
$p_j\ge8\pi^2Q_j-o(Q_j)$ if the marked count identities, curvature
accounting, and all relative boundary terms have total errors
$o(Q_j)$.  Then \eqref{webvar:reduction} gives $\alpha_j/h_j\to4$.
Extra nonnegative energy outside the chosen pieces helps the lower
bound; unaccounted topology does not.  Anderson's theorem supplies
neither these marked whole AE cores nor the required $o(Q_j)$
precision from an $O(A_j)$ budget.  The conditional deduction
therefore leaves the unrestricted O11 coefficient open.
